% !TEX program = xelatex
% 肝癌空间免疫抑制微环境分析项目报告（Nature 风格重构版）
% 使用 ctexart 文档类，XeLaTeX 编译

\documentclass[12pt, a4paper]{ctexart}

% ============================================================
% 宏包加载
% ============================================================
\usepackage{geometry}
\geometry{left=2.5cm, right=2.5cm, top=2.5cm, bottom=2.5cm}

\usepackage{graphicx}
\usepackage{booktabs}
\usepackage{longtable}
\usepackage{amsmath}
\usepackage{amssymb}
\usepackage{hyperref}
\usepackage{xcolor}
\usepackage{listings}
\usepackage{float}
\usepackage{caption}
\usepackage{subcaption}
\usepackage{enumitem}
\usepackage{setspace}
\usepackage{fancyhdr}
\usepackage{natbib}
\usepackage{array}
\usepackage{multirow}
\usepackage{tabularx}
\usepackage{appendix}
\usepackage{microtype}

% ============================================================
% 页面与格式设置
% ============================================================
\onehalfspacing
\pagestyle{fancy}
\fancyhf{}
\fancyhead[L]{\small 基于单细胞与空间转录组的肝癌免疫抑制微环境分析}
\fancyhead[R]{\small \thepage}
\renewcommand{\headrulewidth}{0.4pt}

\lstset{
    basicstyle=\small\ttfamily,
    keywordstyle=\color{blue},
    commentstyle=\color{gray},
    stringstyle=\color{red},
    breaklines=true,
    frame=single,
    numbers=left,
    numberstyle=\tiny\color{gray},
    language=Python,
    showstringspaces=false,
    xleftmargin=1em,
    framexleftmargin=1em,
    extendedchars=false,
    keepspaces=true
}

\hypersetup{
    colorlinks=true,
    linkcolor=blue,
    citecolor=blue,
    urlcolor=blue
}

% ============================================================
% 标题信息
% ============================================================
\title{
    \vspace{-1cm}
    \textbf{基于单细胞与空间转录组整合的肝癌\\
    免疫抑制微环境空间解析与预后验证}\\[0.5em]
    \large 项目技术报告
}
\author{}
\date{\today}

% ============================================================
% 正文开始
% ============================================================
\begin{document}

\maketitle
\thispagestyle{empty}

\begin{abstract}
肝细胞癌（Hepatocellular Carcinoma，HCC）是全球癌症相关死亡的第三大原因，
免疫检查点抑制剂的临床有效率仅约15\%--20\%，
核心瓶颈在于肿瘤微环境中调节性 T 细胞（Treg）、
肿瘤相关巨噬细胞（TAM）与肿瘤相关成纤维细胞（CAF）所共同构建的多层次免疫抑制网络~\citep{llovet2021hepatocellular,finn2020atezolizumab}。
本项目整合单细胞转录组（scRNA-seq）与 10x Genomics Visium 空间转录组数据，
以 Cell2location 贝叶斯反卷积算法~\citep{kleshchevnikov2022cell2location}
将细胞类型注释精确映射至空间切片的每个 spot，
并采用基于 Schürch 等人"细胞邻域（Cellular Neighborhoods）"
框架~\citep{schurch2020coordinated}的邻域组成聚类策略，
在不依赖先验阈值的前提下识别免疫抑制空间生态位。
通过配体-受体通讯定量分析~\citep{efremova2020cellphonedb}
揭示 Treg--TAM--CAF 轴核心信号通路（CXCL12--CXCR4、CCL22--CCR4、TGFB1--TGFBR1、PD-1--PD-L1），
利用两张独立 Visium 切片（HCC4R 与 CHC20）验证跨样本一致性，
并以 ssGSEA 方法将提取的空间特征基因签名投影至 TCGA-LIHC 队列
（约370例患者）~\citep{tomczak2015cancer}进行独立预后验证。
本项目建立了从空间精细解析到机制性通讯分析再到临床队列验证的完整转化分析框架。

\noindent\textbf{关键词：}
肝细胞癌；空间转录组；单细胞转录组；免疫抑制微环境；Cell2location；
Cellular Neighborhoods；配体-受体通讯；ssGSEA；生存分析
\end{abstract}

\newpage
\tableofcontents
\newpage

% ============================================================
%  第一部分：研究背景与意义
% ============================================================

\section{引言}

\subsection{疾病背景与临床挑战}

肝细胞癌（HCC）是全球第六大恶性肿瘤，每年新发病例超过90万，死亡病例约83万，
其中中国约占全球新发病例的45\%~\citep{llovet2021hepatocellular,chen2021liver}。
由于大多数患者确诊时已属晚期，HCC 的五年生存率不足20\%。
索拉非尼曾是晚期 HCC 的标准一线方案，但相比对照组仅将中位总生存期延长约3个月~\citep{llovet2008sorafenib}。
2020年 IMbrave150 试验证明，阿替利珠单抗联合贝伐珠单抗可将客观缓解率提升至27.3\%，
并首次在总生存期上优于索拉非尼~\citep{finn2020atezolizumab}；
然而仍有约70\%的患者对该联合方案无应答，揭示肿瘤微环境中存在强大的多层次免疫抑制机制。

HCC 微环境中的免疫抑制由多种细胞协同构建。调节性 T 细胞（Treg）以 CD$4^+$CD$25^+$ FOXP$3^+$ 为经典表型，通过分泌 IL-10、TGF-$\beta$ 和 IL-35，以及表达 CTLA-4、TIGIT、LAG-3 等检查点分子，
直接抑制效应性 T 细胞与 NK 细胞的抗肿瘤活性~\citep{sakaguchi2020regulatory}；
单细胞转录组研究表明 Treg 在肝癌组织中显著富集，
Treg 与 CD$8^+$ T 细胞的比值是独立预后因子~\citep{zheng2017landscape}。
以 M2 极化表型（CD163$^+$MRC1$^+$）为主的肿瘤相关巨噬细胞（TAM）通过 CCL22--CCR4 和 CXCL12--CXCR4 轴
主动招募 Treg 进入肿瘤区域，同时分泌 TGF-$\beta$ 并高表达 PD-L1 以抑制效应性 T 细胞~\citep{biswas2012macrophage,noy2014tumor}；
TAM 与 Treg 之间的双向正反馈循环共同维持了稳定的局部免疫抑制生态。
以 FAP、ACTA2 和 POSTN 为标志的肿瘤相关成纤维细胞（CAF）则通过分泌 CXCL12 招募免疫抑制细胞，
并通过重塑细胞外基质（ECM）形成物理屏障阻碍效应性 T 细胞浸润~\citep{kalluri2016biology,mariathasan2018tgfb}。
这一 Treg--TAM--CAF 三角轴的协同作用是 HCC 免疫治疗耐药的核心机制，
也是本项目重点解析的生物学对象。

\subsection{空间转录组学与整合分析范式}

上述研究需要理解免疫抑制的空间组织结构。
传统 bulk RNA-seq 将组织均质化后测序，丢失了细胞间的空间位置信息；
scRNA-seq 虽能解析单细胞转录特征，但在组织解离过程中不可避免地破坏了空间关系~\citep{luecken2019current}。
10x Genomics Visium 空间转录组平台在保留组织原位结构的前提下，
对每个空间位点（spot，直径约55\,$\mu$m）进行全转录组测序，
使研究者得以直接观测细胞类型的空间分布模式~\citep{williams2022introduction}。
由于单个 spot 覆盖1--10个细胞的混合信号，
需结合 scRNA-seq 参考数据通过计算反卷积对其细胞组成进行定量估计。

在当前主流反卷积工具中，Cell2location 基于负二项贝叶斯层次模型，
在全基因组范围内同时估计多种细胞类型的绝对丰度，并提供后验不确定性量化~\citep{kleshchevnikov2022cell2location}；
多项独立基准测试一致表明，它在稀有细胞类型（如 Treg，比例通常 $<$5\%）检测上
优于 RCTD~\citep{cable2021robust}、CARD~\citep{ma2022integrative} 和 Tangram~\citep{biancalani2021deep}~\citep{li2022benchmarking,sang2023spotless,yan2023benchmarking}。
Schürch 等人在结直肠癌中发展的"细胞邻域"概念，
通过对局部细胞组成进行无监督聚类揭示不同 neighborhood 对抗肿瘤免疫的差异化贡献，
为空间分辨率下的肿瘤微环境分型提供了方法学范式~\citep{schurch2020coordinated}。
在此基础上，将空间转录组局部特征基因签名通过 ssGSEA~\citep{barbie2009systematic}
投影至大规模临床队列，结合多变量 Cox 回归~\citep{cox1972regression}
进行独立预后验证，是当前实现"空间发现--临床转化"的经典策略~\citep{hanzelmann2013gsva}。

现有 HCC 研究仍普遍缺少三个维度的整合：
在保留空间信息的前提下，将 Treg--TAM--CAF 轴的共定位从描述性层面提升至配体-受体通讯机制层面；
采用数据驱动的方法识别免疫抑制空间生态位；
以及将空间转录组发现系统性地在独立大规模临床队列中进行预后验证。
本项目针对上述不足，建立了一套可推广的完整转化分析框架。

\subsection{研究目标}

本项目旨在回答四个相互关联的核心科学问题：
（1）肝癌 Visium 切片中，免疫抑制细胞信号（Treg、TAM、CAF）是否存在系统性空间共定位，
且富集于特定的肿瘤-基质交界区域？
（2）该区域内 Treg--TAM--CAF 轴是否通过特定配体-受体对（CXCL12--CXCR4、CCL22--CCR4 等）
实现协同免疫抑制？
（3）所识别的免疫抑制空间生态位特征是否在多张独立切片间具有跨样本一致性？
（4）基于空间数据提炼的特征基因签名能否在 TCGA-LIHC 大队列中独立预测患者预后？

% ============================================================
%  第二部分：方法
% ============================================================

\section{方法}

\subsection{整体分析流程}

本研究建立了一套贯穿``单细胞参考构建--空间反卷积--生态位识别--机制解析--跨样本与临床验证"五个阶段的完整转化分析框架（图~\ref{fig:workflow}）。整体流程以 scRNA-seq 参考图谱为起点，经 Cell2location 贝叶斯反卷积将细胞类型丰度映射至 Visium 空间切片的每个 spot；随后基于空间邻域组成的无监督聚类识别免疫抑制生态位，并通过两层次差异表达筛选提炼生态位特征基因签名；在此基础上分别开展配体-受体通讯定量分析（机制层面）、独立切片 CHC20 的跨样本复现验证（空间层面）以及 TCGA-LIHC 大队列的 ssGSEA 投影与多变量 Cox 生存分析（临床层面），三条分析支线相互独立又共同指向同一组生态位特征基因签名，从而构成"发现--机制--验证"的闭环证据链。各阶段对应的具体方法学细节见下文 2.2--2.6 节。

\begin{figure}[h] % h:当前位置, t:顶部, b:底部, p:独立一页。LaTeX 会自动选择最合适的位置
    \centering % 使图片居中
    \includegraphics[width=1.0\textwidth]{plots/workflow_cn.pdf} 
    \caption{整体分析流程图}
    \label{fig:workflow}
\end{figure}

\subsection{数据来源与质量控制}

本研究整合了单细胞转录组（scRNA-seq）、空间转录组（ST）及大规模临床大规模转录组等多模态数据集，构建``外周发现--空间定位--临床验证”的转化医学分析框架。

单细胞参考数据集下载自 GEO 数据库（登录号：GSE149614），涵盖肝细胞（Hepatocyte）、T/NK 细胞、B 细胞、髓系细胞（Myeloid Cell）和成纤维细胞（Fibroblast）等核心细胞类型。参照当前单细胞组学分析的最佳实践~\citep{luecken2019current}，本研究严格过滤了总表达量低于 200 的低质量细胞以及在少于 10 个细胞中表达的基因。鉴于后续所采用的 Cell2location 空间反卷积模型基于贝叶斯负二项回归框架，其内部已实现了对测序深度和技术批次效应的概率建模，因此在单细胞预处理阶段不执行常规的文库大小归一化或 $\log(1+p)$ 变换，直接保留原始 Count 矩阵作为模型的输入。

本研究所使用的空间转录组测序数据均由 10x Genomics Visium 高通量平台生成，采用双中心、跨临床情境的严谨设计构建空间微环境特征的发现与验证闭环。其中，作为主发现集的切片 HCC4R源自前瞻性临床试验队列（GSE238264），该样本来自接受卡博替尼联合纳武利尤单抗靶向免疫新辅助治疗后的HCC切除标本，具有丰富的HCC-CAF空间交互背景；独立外部验证集切片 CHC20源自另一项聚焦于非酒精性脂肪性肝炎相关肝癌的独立空间组学队列，该队列核心关注肿瘤纤维化病灶中骨髓细胞介导的免疫抑制机制。

参照空间转录组分析的最佳实践与质控标准~\citep{williams2022introduction,palla2022squidpy}，对上述两张 ST 切片进行统一过滤：仅保留总 UMI 计数在 5000--35000 之间、且线粒体基因比例低于 20\% 的高质量 Spots，同时剔除在少于 10 个 Spots 中表达的低频基因。

最后，为了将空间组学发现系统性地推向独立大规模临床队列进行生存预后验证，本研究引入了 TCGA-LIHC 队列（包含约 370 例 HCC 患者）。通过 TCGAbiolinks R 包~\citep{colaprico2016tcgabiolinks} 从 GDC 数据库下载标准化的 HTSeq-FPKM 格式表达矩阵以及完整的临床随访与生存信息。

\subsection{Treg 亚群鉴定与 Cell2location 反卷积}

从 scRNA-seq 数据中提取 T/NK 细胞亚群，
以 Leiden 算法（分辨率 res=3.0）进行高分辨率亚聚类~\citep{traag2019leiden}，
通过气泡图可视化 \textit{CD3D}、\textit{CD4}、\textit{FOXP3} 和 \textit{IL2RA} 在各子簇中的表达分布。
确认特定簇呈现 CD4$^+$CD25$^+$FOXP3$^+$ 的 Treg 典型特征后~\citep{sakaguchi2020regulatory}，
将该簇重注释为 Treg 并生成最终细胞类型标签。

Cell2location 分两阶段运行~\citep{kleshchevnikov2022cell2location}。
第一阶段，以 RegressionModel 在 scRNA-seq 数据（原始 count）上训练细胞类型参考签名矩阵，
最大训练轮次250 epochs，batch size 1024，early stopping patience=30；
第二阶段，以参考签名矩阵对每张 Visium 切片独立运行空间建模，
关键先验参数 N\_cells\_per\_location=30（基于 Visium spot 直径约55\,$\mu$m 及典型肝脏组织细胞密度估计~\citep{kleshchevnikov2022cell2location}），
detection\_alpha=20（官方推荐默认值），最大训练1000 epochs。
输出后验均值丰度矩阵（\texttt{means\_cell\_abundance\_w\_sf}）作为各 spot 细胞类型绝对丰度的估计。

为便于结果展示与生物学解读，本研究在 Cellular Neighborhoods 聚类识别免疫抑制生态位之外，进一步以规则化阈值对每个 spot 附加了一个辅助性的空间结构域标签（tumor\_core / tumor\_edge / stroma\_immune / other），该标注仅用于可视化呈现和跨方法交叉验证，不参与 niche 发现主体流程。借鉴Schürch 等人~\citep{schurch2020coordinated} 及 Bhate 等人~\citep{bhate2022tissue} 在多重染色空间组学中采用局部细胞组成阈值区分肿瘤实质、侵袭前沿与基质/免疫区的分域思路，以 Hepatocyte 比例的 75 百分位为阈值，将比例不低于该阈值的 spot 判定为高肝细胞区域（hep\_high）并标注为 tumor\_core；以 immune\_stroma\_score（详见附录~\ref{app:scores}）的第 50 百分位为阈值，在非 hep\_high 的 spot 中，评分不低于该阈值者标注为 stroma\_immune；若该 spot 进一步在物理邻域半径内至少存在一个 hep\_high 邻居，则改判为 tumor\_edge，以刻画肿瘤实质与免疫基质区之间的过渡界面；其余 spot 标注为 other。

\subsection{免疫抑制空间生态位的识别}

以 Cell2location 输出的绝对丰度归一化为比例矩阵后，
利用 $k$ 近邻算法（$k$-NN）在组织几何坐标上构建空间物理邻接图。
对每个 spot $i$ 计算其在物理空间中 $k$ 个最近邻 spot 的细胞类型比例均值，从而抽象出邻域组成向量 $\mathbf{n}_i$：
\begin{equation}
\mathbf{n}_i = \frac{1}{|\mathcal{N}_i|} \sum_{j \in \mathcal{N}_i} \mathbf{p}_j
\label{eq:neighborhood_vector}
\end{equation}
本研究沿用 Schürch 等人~\citep{schurch2020coordinated} 在结直肠癌研究中的经典参数选择 $k=15$。
随后，以该邻域组成向量 $\mathbf{n}_i$ 为高维表型特征，在高维特征空间中重新构建 $k$-NN 图。最后，在此特征邻接图上运行 Leiden 社区检测算法（分辨率$=0.5$）以自适应识别空间生态位类型。
与 KMeans 相比，Leiden 无需预先指定聚类数、且理论上保证了聚类连通性~\citep{traag2019leiden}。

对每个 Leiden 聚类，计算 Treg 比例、Myeloid 比例、
Fibroblast 比例和免疫抑制基因模块评分四个功能维度的均值，
按各维度降序排名后求和得到综合排名 $R_c$：
\begin{equation}
R_c = \text{rank}(\bar{p}_{\text{Treg},c}) + \text{rank}(\bar{p}_{\text{Myeloid},c})
+ \text{rank}(\bar{p}_{\text{Fibroblast},c}) + \text{rank}(\bar{s}_{\text{immune},c})
\label{eq:combined_rank}
\end{equation}
$R_c$ 最小的 cluster 注释为免疫抑制生态位，该多维综合评分确保所识别的 niche 同时满足 Treg--TAM--CAF 三角轴的协同富集特征，
而非单一细胞类型的孤立高表达。

免疫抑制基因模块评分由14个文献支持的标志基因（\textit{FOXP3}、\textit{IL2RA}、\textit{CTLA4}、\textit{TIGIT}、\textit{LAG3}、\textit{PDCD1}、\textit{HAVCR2}、\textit{TGFB1}、\textit{IL10}、\textit{CXCL12}、\textit{CCL22}、\textit{TNFRSF18}、\textit{TNFRSF4}、\textit{IKZF2}）~\citep{sakaguchi2020regulatory,sharma2023immunosuppressive}
的 CP10K+log1p 表达均值定义；
辅助评分体系（Treg-like score、immune-stroma score 和 niche score）通过对相应组成维度进行 Z-score 标准化后加和计算（详见附录）。

为验证上述邻域参数选择不会主导结论，本研究设计了不依赖硬阈值截断的连续型敏感性分析。分别以 $k$-NN 邻居数 $k \in \{10, 15, 20\}$ 及半径倍增系数 $\text{radius\_multiplier} \in \{1.0, 1.25, 1.5\}$ 两种邻域定义方式对综合 niche 评分（公式~\ref{eq:combined_rank} 所依据的四维 Z-score 加总）进行参数扰动，以主分析参数（$k=15$）下的评分向量作为参考基准，采用 Spearman 秩相关系数 $\rho$ 与连续型加权 Jaccard 相似度 WJ 两个互补指标（公式定义见附录~\ref{app:sensitivity}）联合量化每次扰动后评分分布与参考基准的一致性。两指标联合解读的稳健性判定标准为：$\rho > 0.95$ 且 $\text{WJ} > 0.90$ 判定为结果极稳健；$\rho > 0.90$ 且 $\text{WJ} > 0.80$ 判定为结果稳健、参数选择不敏感；$\rho < 0.85$ 或 $\text{WJ} < 0.70$ 则判定为该参数变化对结果有实质性影响，需要重新审视参数选择~\citep{spearman1904proof}。本研究以 $\rho > 0.90$ 且 $\text{WJ} > 0.80$ 作为参数稳健性的通过标准，在 $k$-NN 和半径两种邻域定义、共计 6 组参数扰动下的评分向量与主分析参考基准逐一比较，用于判断 $k=15$ 的参数选择是否显著影响下游 niche 识别与特征基因提取结果。

在完成邻域参数稳健性验证的基础上，本研究进一步以连续型综合 niche 评分（immunosuppressive\_niche\_score）为依据，对每个 spot 附加二值化的高/低生态位标签，作为后续特征基因两层筛选（2.6 节）的分组变量。具体而言，以该评分在全部 spot 中的第 90 百分位为阈值，评分不低于该阈值的 spot 标注为 niche\_high，其余标注为 niche\_low。

\subsection{特征基因签名的两层筛选}

Visium 每个 spot 覆盖约 5--50 个细胞的混合信号（细胞稀释效应），
使 \textit{FOXP3}、\textit{CTLA4} 等 Treg 标志基因在绝大多数 spot 中的原始表达量为零，
导致两组间均值差异（$\log_2$FC）被严重稀释，传统仅依赖 Wilcoxon 检验 + $\log_2$FC 阈值的
差异表达筛选策略难以将这些低丰度但生物学关键的基因纳入 Top-N 候选列表。
为此，本研究在 $\log_2$FC 之外引入检出率（fraction of spots with expression $>0$）
作为互补维度，并以两层筛选策略提取免疫抑制生态位的特征基因签名。

第一层（Wilcoxon + BH-FDR + AUC + 检出率综合排序）：
首先以空间网格聚合（block\_mwu 模式）将地理位置相邻的连续 spot 合并为 block 级别均值，
以降低空间自相关对假设检验独立性的破坏~\citep{palla2022squidpy}；
基于免疫抑制生态位综合评分（公式~\ref{eq:combined_rank} 所依据的四维 Z-score 加总）
的第 90 百分位阈值，将 block 划分为
niche\_high（$\text{score} \geq P_{90}$）与 niche\_low（$\text{score} < P_{90}$）两组，
对每个基因执行单侧 Mann-Whitney U 检验（即 Wilcoxon 秩和检验），
并以 Benjamini-Hochberg 方法对全部基因的 $p$ 值进行 FDR 多重检验校正~\citep{mann1947test}；
在此基础上，分别计算两组的检出率 $\text{frac}_{\text{high}}$、$\text{frac}_{\text{low}}$
及其差值 $\Delta\text{frac} = \text{frac}_{\text{high}} - \text{frac}_{\text{low}}$，
以及由 Mann-Whitney U 统计量标准化得到的 AUC~\citep{hanley1982meaning}；
主筛选条件为 FDR $<$ 0.1 且 AUC $>$ 0.60 且（$\log_2$FC $>$ 0.3 OR $\Delta$frac $>$ 0.10），
若满足条件的基因数不足 80 个，则放宽至 FDR $<$ 0.2 且 AUC $>$ 0.55 且
（$\log_2$FC $>$ 0.2 OR $\Delta$frac $>$ 0.05）；
通过筛选的候选基因进一步剔除涵盖管家基因、核糖体蛋白、线粒体基因、血红蛋白基因、
细胞周期基因、广谱炎症基因及组织损伤/应激基因等350余个非特异性基因的黑名单，
以确保签名基因的 niche 特异性而非组织普遍活跃性；
最终以
\begin{equation*}
\text{signature\_rank\_score} = 0.35 \cdot \text{minmax}(\log_2\text{FC}) + 0.35 \cdot \text{minmax}(\Delta\text{frac}) + 0.30 \cdot \text{minmax}(\text{AUC})
\label{eq:signature_rank_score}
\end{equation*}
对候选基因进行加权综合排序，取排序前 80 个基因（若强制纳入的先验基因导致总数超出，则不作截断）作为初步签名基因集。

第二层（先验功能基因集的靶向 AUC 检验与强制补全）：
针对 Treg 标志物（\textit{FOXP3}、\textit{IL2RA}、\textit{CTLA4} 等）、TAM 特征基因（\textit{CD163}、\textit{MRC1}、\textit{TGFB1} 等）
和 CAF 激活基因（\textit{FAP}、\textit{ACTA2}、\textit{COL1A1} 等）三组具有明确文献支持的先验功能基因集，
在 niche\_high 与 niche\_low 两组间逐一计算 AUC、$\log_2$FC、检出率差值 $\Delta$frac 及
FDR 校正后的 $p$ 值，用于独立评估每个先验基因是否体现出 niche 特异性富集；
对满足 AUC $>$ 0.60 且 $\Delta$frac $>$ 0.05（且不属于非特异性基因黑名单）的先验基因，
无论其是否已通过第一层的全基因组排序，均强制并入最终签名基因列表，
以防止因 Visium 稀疏表达导致部分生物学意义明确的关键基因被漏选。
两层筛选流程的完整实现见附录~\ref{app:code}。

\subsection{配体-受体通讯分析}

为将上述描述性的细胞共定位结论提升至机制性通讯层面，本研究针对 \textit{CXCL12}--\textit{CXCR4}、\textit{CCL22}--\textit{CCR4}、\textit{TGFB1}--\textit{TGFBR1}、PD-1--PD-L1 等 12 对具有明确文献支持的配体-受体（L-R）通讯对（完整列表见附录~\ref{app:lr_table}），在 niche\_high 与 niche\_low 两组间定量比较其通讯强度。考虑到配体-受体乘积法~\citep{efremova2020cellphonedb} 在 Visium spot 级别数据中易因任一基因表达接近零（细胞稀释效应）而使乘积恒为零、导致 \textit{CCL22}--\textit{CCR4} 等 Treg 特异性通路信号完全湮灭，本研究对每个基因的表达向量先在全部 spot 范围内做秩归一化（取值映射至 $[0,1]$），再以配体秩与受体秩之积作为通讯强度的代理指标：
\begin{equation*}
\text{score}_{LR,i} = \text{rank}(E_{L,i}) \times \text{rank}(E_{R,i}), \quad \text{rank}(\cdot) \in [0,1]
\end{equation*}
该秩归一化乘积在保留配体与受体表达相对高低顺序信息的同时，有效规避了稀疏零值对乘积的支配效应，并天然满足配体、受体须同时具有较高表达水平这一逻辑 AND 约束。对每个 L-R 对分别计算其在 niche\_high 和 niche\_low 两组中的秩归一化乘积均值，并以 $\log_2(\text{mean}_{\text{high}}/\text{mean}_{\text{low}})$（分母加入伪计数 $\varepsilon = 10^{-6}$ 以避免除零）量化其在免疫抑制生态位中的富集或耗竭方向，最终以热图（通讯强度）与条形图（$\log_2$FC）联合呈现 12 对通讯轴在两组间的差异模式。

在细胞间通讯之外，本研究进一步以预定义功能基因集评分的方式，从通路层面系统刻画免疫抑制、基质重塑、炎症与代谢应激等生物学过程在不同空间结构域中的活跃程度（图~\ref{fig:domain_func}）。具体而言，参照 Hallmark 基因集~\citep{liberzon2015molecular} 的构建思路及单样本基因集评分的通用范式~\citep{barbie2009systematic}，人工组织了 TGF-$\beta$ Signaling、EMT、Angiogenesis、IL6-JAK-STAT3、Treg Immune Suppression、TAM Polarization (M2)、CAF Activation、Immune Checkpoint、Hypoxia 及 TNF-$\alpha$ Signaling via NF-$\kappa$B 共 10 条与 Treg--TAM--CAF 免疫抑制轴及其伴随微环境改变密切相关的功能通路基因集，每条通路包含 6 个代表性基因；对每个 spot 计算通路内全部可用基因的 CP10K+log1p 表达均值作为该通路的单样本评分，再按 2.5 节所述规则化空间域标签（tumor\_core / tumor\_edge / stroma\_immune / other）分组求平均，得到每条通路在四个空间域中的平均评分。气泡图中气泡大小表征评分的相对高低，颜色区分该评分是否高于该通路在全部四域中的全局均值（红色为高于均值、蓝色为低于均值）。

\subsection{跨切片一致性验证与 TCGA 生存分析}

以 HCC4R 提取的签名基因集为基准，将免疫模块基因（Treg/TAM 相关，共21个）
和基质/ECM 模块基因（CAF/ECM 相关，共21个）分别在 CHC20 切片上计算模块评分，
以 Z-score 标准化后加和得到组合 niche 评分（immune\_stromal\_niche\_score），
通过空间投影图和 Mann-Whitney U 检验~\citep{mann1947test}
量化评分分布的跨样本一致性。

生存分析使用 GSVA R 包~\citep{hanzelmann2013gsva} 的 ssGSEA 方法~\citep{barbie2009systematic}，
以 $\log_2(\text{FPKM}+1)$ 矩阵为输入，高斯核（kcdf=Gaussian）、绝对排名（abs.ranking=TRUE）
对 TCGA-LIHC 队列中每位患者计算标量富集评分。
ssGSEA 为每位患者独立生成一个反映生态位活性的连续评分，
更适合纳入预后回归模型，而非 GSVA 的跨样本相对活性估计~\citep{barbie2009systematic}。
构建多变量 Cox 比例风险回归模型~\citep{cox1972regression}，
以 ssGSEA 评分、年龄和肿瘤分期为协变量：
\begin{equation*}
h(t \mid \mathbf{x}) = h_0(t) \cdot \exp(\beta_1 s_{\text{ssGSEA}} + \beta_2 \text{age} + \beta_3 \text{stage})
\label{eq:cox}
\end{equation*}
控制年龄和分期（HCC 已知强效独立预后因子~\citep{llovet2021hepatocellular}）
后，若 ssGSEA 评分的 $\beta_1$ 仍显著，则证明免疫抑制生态位签名捕获了独立于临床分期以外的预后信息。
以 ssGSEA 评分中位数为阈值，绘制 Kaplan-Meier 生存曲线~\citep{kaplan1958nonparametric}，
log-rank 检验~\citep{mantel1966evaluation} 评估两组生存差异。

% ============================================================
%  第三部分：结果
% ============================================================

\section{结果}

\subsection{scRNA-seq 参考图谱构建与 Treg 亚群鉴定}

\begin{figure}[h]
    \centering
    % 第一张子图
    \begin{subfigure}[b]{0.48\textwidth}
        \centering
        \includegraphics[width=\textwidth]{plots/cell_type_marker_gene_expression.png} % 填入第一张图片名
        \caption{主要细胞类型的特征标志基因表达热图} % 子图 (a) 的标题
        \label{fig:cell_type} % 子图 (a) 的标签
    \end{subfigure}
    \hfill % 关键：利用弹性间距将两张图推向两端，实现横向并列
    % 第二张子图
    \begin{subfigure}[b]{0.48\textwidth}
        \centering
        \includegraphics[width=\textwidth]{plots/treg_marker_gene_expression_t_nk_subclusters.png} % 填入第二张图片名
        \caption{T/NK 细胞亚群中 Treg 标志基因表达气泡图} % 子图 (b) 的标题
        \label{fig:treg} % 子图 (b) 的标签
    \end{subfigure}
    \caption{单细胞参考图谱的细胞类型鉴定与调节性 T 细胞（Treg）亚群的精细识别}
    \label{fig:scRNA-seq_Treg}
\end{figure}

利用 GSE149614 单细胞转录组数据构建细胞类型参考图谱。
细胞类型 Marker 基因热图（图 \ref{fig:cell_type}）以 Z-score 形式展示各细胞类型的平均表达特征，横轴为识别出的 7 种主要细胞类型（包括 T/NK、Treg、Myeloid 等），纵轴为各类型经典的 Marker 基因，颜色代表经 Z-score 归一化后的相对表达水平。其中，肝细胞/恶性细胞（\textit{ALB}、\textit{APOA2}、\textit{APOA1}）、T/NK 细胞（\textit{CD3D}、\textit{GNLY}、\textit{GZMB}）、髓系细胞（\textit{LYZ}、\textit{AIF1}、\textit{HLA-DRA}）、HSC/Fibroblast（\textit{COL1A1}、\textit{ACTA2}、\textit{RGS5}）、B 细胞（\textit{MS4A1}、\textit{CD79A}）和内皮细胞（\textit{PECAM1}、\textit{CDH5}）均呈现清晰的对角线高亮模式，验证了细胞类型注释的可靠性~\citep{ma2021comprehensive}。

针对 T/NK 亚群，以 res=3.0 进行高分辨率 Leiden 聚类并绘制横版气泡图（图 \ref{fig:treg}），展示了不同 T/NK 亚群（Subclusters）中 \textit{CD3D}、\textit{CD4}、\textit{FOXP3} 和 \textit{IL2RA} 的表达情况。其中，气泡大小代表细胞表达比例（Fraction of cells），颜色深浅代表平均表达量。特定簇（簇9）呈现 \textit{FOXP3} 气泡最大且颜色最深、\textit{CD4} 共表达、\textit{IL2RA} 显著高于其他簇的组合模式，符合 CD4$^+$CD25$^+$FOXP3$^+$ 经典 Treg 表型~\citep{sakaguchi2020regulatory}；其余 T/NK 亚簇 \textit{FOXP3} 近乎缺失，证明了 \textit{FOXP3} 表达的亚群特异性。将簇9重注释为 Treg，生成最终细胞类型标签，该结果纳入 Cell2location 反卷积的参考签名中。

\subsection{Cell2location 反卷积揭示 HCC4R 的免疫细胞空间分布}
\subsubsection{HCC4R 切片免疫细胞的空间分布}
\begin{figure}[h] % h:当前位置, t:顶部, b:底部, p:独立一页。LaTeX 会自动选择最合适的位置
    \centering % 使图片居中
    \includegraphics[width=0.6\textwidth]{plots/hcc4r_he_image.png} 
    \caption{肝细胞癌（HCC4R）空间生态位与肿瘤微环境的解剖轮廓划分}
    \label{fig:he}
\end{figure}

图 \ref{fig:he} 将 HCC4R 切片中的空间生态位地图精确映射每个 spot 上。其中的 H\&E 染色图清晰呈现了肿瘤实质区（核大深染、Hepatocyte 密集）与基质/间质区（纤维化、炎症浸润）的形态学边界，为空间转录组结果提供了组织学参照。

从空间结构域地图可见，Tumor Core（红色）呈大片连续分布，构成组织切片的主体骨架，与肿瘤实质在 H\&E 染色和 Hepatocyte 反卷积丰度图中呈现的核密集区域在空间位置上高度吻合；Stroma/Immune（蓝色）则以条带状镶嵌于 Tumor Core 周围及组织边缘，形成环绕/穿插于肿瘤实质之间的连续区域，未表现为孤立的散点分布，这与无监督聚类能够有效捕获组织连续空间微结构、而非随机噪声的方法学预期一致~\citep{schurch2020coordinated}；Tumor Edge（橙色）散布于 Tumor Core 与 Stroma/Immune 的过渡地带，构成两者之间的窄带状界面，在图中占比相对最小。这一空间格局提示，HCC4R 切片的组织结构并非肿瘤与基质截然二分，而是呈现"核心—边缘—基质免疫"的连续过渡梯度，与经典``免疫排斥型"HCC 中肿瘤细胞占据核心、免疫基质细胞被限制在外周边界的空间结构描述相符~\citep{llovet2021hepatocellular,binnewies2018understanding}。

\begin{figure}[h]
    \centering
    % === 第一行：Treg 和 Myeloid 细胞 ===
    \begin{subfigure}[b]{0.48\textwidth}
        \centering
        \includegraphics[width=\textwidth]{plots/spatial_hepatocyte.png}
        \caption{Hepatocyte 细胞空间比例分布}
        \label{fig:spatial_hepatocyte}
    \end{subfigure}
    \hfill
    \begin{subfigure}[b]{0.48\textwidth}
        \centering
        \includegraphics[width=\textwidth]{plots/spatial_treg.png}
        \caption{Treg-like 细胞空间丰度分布}
        \label{fig:spatial_treg}
    \end{subfigure}

    \vspace{0.3cm} % 上下两行图片之间的垂直间距

    % === 第二行：Fibroblast 和 Hepatocyte 细胞 ===
    \begin{subfigure}[b]{0.48\textwidth}
        \centering
        \includegraphics[width=\textwidth]{plots/spatial_myeloid.png}
        \caption{Myeloid 细胞空间比例分布}
        \label{fig:spatial_myeloid}
    \end{subfigure}
    \hfill
    \begin{subfigure}[b]{0.48\textwidth}
        \centering
        \includegraphics[width=\textwidth]{plots/spatial_fibroblast.png}
        \caption{Fibroblast 细胞空间比例分布}
        \label{fig:spatial_fibroblast}
    \end{subfigure}
    % 大图总题注
    \caption{基于贝叶斯反卷积算法映射的主要细胞类型空间丰度及比例分布图}
    \label{fig:spatial_cell_abundance}
\end{figure}

细胞类型空间分布图（图 \ref{fig:spatial_cell_abundance}）揭示了与预期高度一致的空间格局。
Hepatocyte 高丰度区域集中于肿瘤实质，与 H\&E 图中核密集区域的空间吻合（图 \ref{fig:spatial_hepatocyte}）；
Treg 呈局灶性富集分布（右偏分布），高丰度热点主要定位于肿瘤实质与间质的交界区域（图 \ref{fig:spatial_treg}），
与 Treg 优先聚集于``肿瘤-免疫界面"的生物学预期一致~\citep{zheng2017landscape}；
Myeloid 细胞与 Treg 的空间分布具有明显共定位趋势（图 \ref{fig:spatial_myeloid}），呼应了 TAM 通过 CCL22--CCR4 轴主动招募 Treg 的信号机制~\citep{noy2014tumor}；
Fibroblast 则主要分布于围绕肿瘤实质的基质层（图 \ref{fig:spatial_fibroblast}），反映 CAF 构建物理屏障的功能定位~\citep{kalluri2016biology}。

\begin{figure}[h] % h:当前位置, t:顶部, b:底部, p:独立一页。LaTeX 会自动选择最合适的位置
    \centering % 使图片居中
    \includegraphics[width=0.8\textwidth]{plots/cell_type_spatial_colocalization.png} 
    \caption{空间组织中主要细胞类型的空间共定位相关性热图}
    \label{fig:hot_plot}
\end{figure}

细胞类型间的 Pearson 相关分析（图 \ref{fig:hot_plot}）进一步以定量方式刻画了上述空间共定位关系。相关性热图结果显示，Myeloid 细胞分别与 Treg 细胞（$r = 0.35$）及 Fibroblast（$r = 0.35$）均呈现出中等强度的显著正相关，而 Fibroblast 与 Treg 之间则表现出微弱的负相关（$r = -0.15$）。这一"非对称三角"相关模式提示，免疫基质组分在空间上可能存在精细的功能分层：Myeloid 细胞（TAM）作为核心的空间桥梁同时与免疫抑制细胞和基质细胞发生关联——一方面通过分泌 CCL22、CXCL12 等趋化因子招募并锚定 Treg~\citep{noy2014tumor,curiel2004specific}，另一方面与 Fibroblast（CAF）协同参与细胞外基质重塑与物理屏障构建~\citep{kalluri2016biology}；而 Fibroblast 与 Treg 二者直接的空间关联较弱，与二者更可能通过 TAM 这一共同节点间接耦联、而非直接相互募集的机制假设相符，这与既往研究中 TAM 在肿瘤免疫抑制网络里承担"信号枢纽"角色的观察一致~\citep{biswas2012macrophage}。

此外，上述三类免疫基质组分（Myeloid、Treg、Fibroblast）均与实质细胞 Hepatocyte 表现出强烈的空间负相关（相关系数分别为 $r = -0.73$、$r = -0.62$ 和 $r = -0.55$），提示三者在空间上高度聚集于非肿瘤实质区域而系统性地排斥于肿瘤核心之外。这一定量证据有力支持了免疫抑制生态位主要盘踞于肿瘤-基质过渡带、而非弥散分布于整个切片的空间格局，其阻碍效应性免疫细胞向肿瘤核心有效浸润的拓扑学特征，与经典``免疫排斥型"肿瘤中免疫细胞被物理限制于肿瘤边界及基质区、无法进入肿瘤实质发挥杀伤功能的空间结构描述高度吻合~\citep{llovet2021hepatocellular,chen2017oncology}。

\subsubsection{HCC4R 不同空间结构域差异}

\begin{figure}[h]
    \centering
    % 第一张图
    \begin{subfigure}[b]{0.8\textwidth}
        \centering
        \includegraphics[width=\textwidth]{plots/cell_type_composition_per_spatial_domain.png}
        \caption{空间结构域的细胞类型组成}
        \label{fig:cell_type_composition}
    \end{subfigure}
    \vspace{0.5cm}
    
    \begin{subfigure}[b]{0.32\textwidth}
        \centering
        \includegraphics[width=\textwidth]{plots/treg_like_score_by_spatial_domain.png}
        \caption{Treg-like 得分}
        \label{fig:treg_score}
    \end{subfigure}
    \hfill % 自动填满中间的空白，使子图两端对齐
    % 第二张图
    \begin{subfigure}[b]{0.32\textwidth}
        \centering
        \includegraphics[width=\textwidth]{plots/immune_stroma_score_by_spatial_domain.png}
        \caption{免疫-基质得分}
        \label{fig:immune_stroma}
    \end{subfigure}
    \hfill % 自动填满中间的空白
    % 第三张图
    \begin{subfigure}[b]{0.32\textwidth}
        \centering
        \includegraphics[width=\textwidth]{plots/immunosuppressive_niche_score_by_spatial_domain.png}
        \caption{免疫抑制微环境得分}
        \label{fig:immunosuppressive_niche}
    \end{subfigure}
    \caption{不同空间结构域的细胞类型组成和不同模块得分}
    \label{fig:spatial_domain_scores}
\end{figure}

细胞组成堆叠条形图（图 \ref{fig:cell_type_composition}）为上述空间格局提供了定量支持。其中四个域中 Hepatocyte 比例呈现明确的梯度递减：Tumor Core（约 35\%）> Other（约 24\%）> Tumor Edge（约 25\%）> Stroma/Immune（最低），与``Tumor Core 由肝细胞/恶性细胞主导"的定义相符。与之相对，Treg、Myeloid 和 Fibroblast 三类免疫抑制轴核心细胞的比例均在 Stroma/Immune 和 Tumor Edge 两个域中高于 Tumor Core：Treg 比例在 Stroma/Immune 中达到四组最高水平，Myeloid 和 Fibroblast 亦呈现相似的富集趋势，三者在同一空间域中的协同富集直接印证了正文所述 Treg--TAM--CAF 三角免疫抑制轴的空间协同性~\citep{sakaguchi2020regulatory,noy2014tumor,kalluri2016biology}。T/NK 细胞比例在 Tumor Edge 中相对突出，与肿瘤侵袭前沿是效应 T 细胞与调节性细胞交锋的主战场这一认识相符~\citep{schurch2020coordinated}。Endothelial 比例在四组间相对稳定，提示血管分布本身并非驱动该聚类分区的主导因素。

空间结构域在不同模块得分小提琴图（图 \ref{fig:treg_score} \ref{fig:immune_stroma} \ref{fig:immunosuppressive_niche}）也可以得到类似的结论。对不同空间结构域的定量比较表明，各区域的微环境特征存在显著差异。Kruskal-Wallis 检验结果显示，三种评分在不同区域间的整体分布差异均具有极显著的统计学意义（均 $p < 0.001$）。其中，Stroma/Immune 区域的免疫-基质评分（$\eta^2 = 0.83$）、免疫抑制微环境评分（$\eta^2 = 0.41$）以及 Treg 样细胞评分（$\eta^2 = 0.33$）均显著高于 Tumor Core、Tumor Edge 及 Other 区域（两两比较均 $p < 0.001$）。这种差异在免疫-基质评分中最为突出：非 Stroma/Immune 区域高度集中于低分区间，而 Stroma/Immune 区域则呈现出明显的右尾偏移，广泛分布于高分区间。~\citep{mann1947test}。

综合细胞组成与功能评分两方面的定量结果，HCC4R 切片呈现出"Tumor Core 低、Tumor Edge 居中、Stroma/Immune 高"的单调免疫抑制梯度，且该规则化分域方案在效应量与显著性上均获统计学支持，为下一节数据驱动的 Cellular Neighborhoods 聚类提供了先验预期。

\subsection{Cellular Neighborhoods 聚类识别免疫抑制空间生态位}

\begin{figure}[h] % h:当前位置, t:顶部, b:底部, p:独立一页。LaTeX 会自动选择最合适的位置
    \centering % 使图片居中
    \includegraphics[width=0.8\textwidth]{plots/neighborhood_composition_clusters.png} 
    \caption{免疫抑制空间生态位 Leiden 聚类结果}
    \label{fig:niche_cluster}
\end{figure}

图 \ref{fig:niche_cluster} 展示了 2.4 节邻域组成向量 $\mathbf{n}_i$（$k$-NN，$k=15$）在特征空间中重新构图后运行 Leiden 聚类（分辨率 $= 0.5$）得到的聚类结果。本图直接以聚类编号（NC-0 至 NC-11，共 12 个 neighborhood cluster）在空间坐标系中着色展示，完整地保留了 Leiden 聚类算法在邻域组成特征空间中原生发现的组织异质性结构。

图中最突出的 Immunosuppressive Niche 是在空间上面积最大、连通性最强的单一聚类区域，占据切片中心地带并向多个方向延伸，与其余 11 个聚类相比在空间连续性和覆盖范围上均具有明显优势。该聚类被正文方法学流程通过综合排名公式（公式 \ref{eq:combined_rank}）判定为 immunosuppressive\_niche ，即在 Treg 比例、Myeloid 比例、Fibroblast 比例和免疫抑制基因评分四个维度的降序排名之和最小，因此在四维证据的联合权衡下被判定为综合免疫抑制特征最强的邻域类型。Immunosuppressive Niche 在空间上包含于图 \ref{fig:he}  Stroma/Immune 和 Tumor Edge 区域中，两种独立标注方式（Leiden 聚类原始编号 vs. 规则化语义标签）在空间位置上的相互印证，从方法学角度增强了该免疫抑制生态位判定结果的稳健性，也与聚类图与语义标注图应呈现高度空间一致性的预期相符~\citep{schurch2020coordinated,goltsev2018deep}。

图中同时清晰呈现的一个特征是，除 NC-2 外，其余聚类呈现出明显的组织学分区意义：例如 NC-1 和 NC-3 分别在切片左下与右上角形成两个独立的大面积连续区域，很可能对应不同批次或不同区域的肿瘤实质亚型；NC-5、NC-9等则以较窄条带形式镶嵌于 NC-1、NC-2 等大聚类之间，提示这些邻域代表了肿瘤与基质之间的过渡带亚结构。整体而言，Leiden 算法在未设定聚类数目先验的情况下自动识别出 12 个具有内部一致性、且在空间上连续分布的邻域类型，这一结果验证了无监督聚类方法确实能够捕获真实存在的组织空间微结构而非人为设定的伪影~\citep{traag2019leiden,schurch2020coordinated}。这种``一个显著免疫抑制邻域 + 多个肿瘤/基质亚结构邻域"并存的格局，也与近年空间蛋白质组学研究中发现的肿瘤组织普遍存在多种细胞邻域共存、且免疫抑制性邻域通常表现为跨越肿瘤-基质界面的连续结构而非孤立热点这一观察一致~\citep{jackson2020single,bhate2022tissue}，进一步支持了本研究识别出的免疫抑制生态位是 HCC4R 切片中一种具有稳定组织学基础的真实空间实体，而非聚类参数任意选择下的偶然产物。

\subsection{niche 特征基因签名的提取与功能注释}

\begin{figure}[h]
    \centering
    \begin{subfigure}[b]{0.71\textwidth}
        \centering
        \includegraphics[width=\textwidth]{plots/fraction_of_spots_enrichment.png}
        \caption{高低生态位特征基因筛选散点图}
        \label{fig:high_low_niche_gene}
    \end{subfigure}
    \hfill % 自动填满中间的空白，使子图两端对齐
    \begin{subfigure}[b]{0.27\textwidth}
        \centering
        \includegraphics[width=\textwidth]{plots/niche_signature_genes_by_spatial_domain.png}
        \caption{核心差异基因在各空间域的表达丰度热图}
        \label{fig:domain_hotplot}
    \end{subfigure}
    \caption{肿瘤微环境不同微结构区域的差异表达基因识别与景观图}
    \label{fig:niche_core_gene}
\end{figure}

图 \ref{fig:high_low_niche_gene} 以 $\log_2$FC 为横轴、Δfraction（检出率差值）为纵轴，展示了 Layer 1 检出率补偿维度的筛选结果：红色点（Both-sig，$n = 33$）同时满足双重显著条件，包括 \textit{CXCL12}、\textit{TGFB1} 等高置信免疫抑制基因，与已知趋化因子轴和 CAF/Treg 分泌机制一致~\citep{domanska2013targeting,mariathasan2018tgfb}。值得注意的是，多数免疫抑制候选基因的 $\log_2$FC 仅集中在 $0.1$--$0.3$ 之间但 Δfraction 已较高，这一反常可由 Visium 的细胞稀释效应解释~\citep{luecken2019current}，说明 Δfrac 维度对传统 $\log_2$FC 筛选具有必要的补偿作用。

图 \ref{fig:domain_hotplot} 展示了按 composite\_score 降序排列的 Top 基因在四个空间域中的表达 Z-score 热图，几乎所有签名基因均呈现 Stroma/Immune 最高、Tumor Core 最低、Tumor Edge 居中的一致趋势，与图 \ref{fig:he} 的空间格局高度吻合~\citep{zheng2017landscape}。热图同时显示 MHC-II 类分子、淋巴细胞相关基因（\textit{CD52}、\textit{TRAC} 等）与基质/成纤维细胞基因（\textit{DCN}、\textit{LUM} 等）在 Stroma/Immune 域共富集，提示该区域兼具抗原呈递活性与基质重塑特征~\citep{neefjes2011towards,iozzo2015proteoglycan}。

\begin{figure}[h] % h:当前位置, t:顶部, b:底部, p:独立一页。LaTeX 会自动选择最合适的位置
    \centering % 使图片居中
    \includegraphics[width=0.9\textwidth]{plots/pathway_enrichment_by_spatial_domain.png} 
    \caption{各空间域中预定义功能通路基因集评分的富集图谱}
    \label{fig:domain_func}
\end{figure}

图 \ref{fig:domain_func} 以气泡图展示了 10 条预定义功能通路基因集评分在四个空间域中的富集情况。全局来看，Tumor Core 在全部 10 条通路中均为最低活跃度，而除 Hypoxia 外的 9 条通路均在 Stroma/Immune 域达到峰值，印证该区域是促癌、基质重塑及免疫抑制信号最核心的活跃场所。在分化模式上，CAF Activation 仅在 Stroma/Immune 域富集，与 CAF 分泌 CXCL12、重塑 ECM 的机制一致~\citep{kalluri2016biology}；EMT、TGF-$\beta$、Treg Immune Suppression 等通路在 Tumor Edge 与 Stroma/Immune 均偏高，符合侵袭前沿与免疫抑制核心区协同活跃的预期~\citep{nieto2016emt,sakaguchi2020regulatory}。值得注意的是 Hypoxia 通路反而在 Tumor Edge 富集而非 Stroma/Immune，可能与该区域血管新生活跃、氧耗-血供矛盾尖锐更易检测到急性 HIF 通路激活有关~\citep{palazon2017hif}。

综上，本节提取的基因签名与功能通路评分共同证实了 Stroma/Immune 作为免疫抑制核心场所、Tumor Edge 作为侵袭前沿过渡带的空间定位，为后续跨样本验证与生存分析提供了严格降噪筛选的特征基因集~\citep{zheng2017landscape,schurch2020coordinated}。

\subsection{niche gene program 在独立切片 CHC20 中的空间复现}

本节展示了 2.7 节中跨切片一致性验证的结果，检验以 HCC4R 切片提取的免疫抑制生态位签名基因集是否具有超越单一切片的跨样本普适性。

\begin{figure}[h] % [htbp] 允许 LaTeX 自动选择最佳位置（当前位置、顶部、底部或单独一页）
    \centering
    % 第一行左图 (a)
    \begin{subfigure}[b]{0.48\textwidth}
        \centering
        \includegraphics[width=\textwidth]{plots/chc20_immune_module_spatial_score.png}
        \caption{免疫模块评分}
        \label{fig:immune_score}
    \end{subfigure}
    \hfill % 水平填充，自动推开左右两张图
    % 第一行右图 (b)
    \begin{subfigure}[b]{0.48\textwidth}
        \centering
        \includegraphics[width=\textwidth]{plots/chc20_stromal_ecm_module_spatial_score.png}
        \caption{基质 ECM 模块评分}
        \label{fig:stromal_score}
    \end{subfigure}
    \vspace{0.4cm} % 上下两行图片之间的垂直间距，可根据需要微调

    % 第二行左图 (c)
    \begin{subfigure}[b]{0.48\textwidth}
        \centering
        \includegraphics[width=\textwidth]{plots/chc20_combined_immune_stromal_niche_score.png}
        \caption{组合微环境评分}
        \label{fig:combined_score}
    \end{subfigure}
    \hfill
    % 第二行右图 (d)
    \begin{subfigure}[b]{0.48\textwidth}
        \centering
        \includegraphics[width=\textwidth]{plots/hcc4r_vs_chc20_combined_niche_score_distribution.png}
        \caption{评分分布对比}
        \label{fig:score_comparison}
    \end{subfigure}

    % 整体中文大标题
    \caption{微环境模块评分的空间分布与对比}
    \label{fig:main_microenvironment_figure}
\end{figure}

图 \ref{fig:immune_score}、\ref{fig:stromal_score} 和 \ref{fig:combined_score} 分别展示了将签名基因按照功能来源拆分后的免疫模块评分（$n = 15$ 个匹配基因，涵盖 \textit{FOXP3}、\textit{IL2RA}、\textit{CTLA4}、\textit{TIGIT}、\textit{PDCD1}、\textit{LAG3}、\textit{CD163}、\textit{MRC1}、\textit{CXCL12} 等 Treg/TAM 核心分子）、基质/ECM 模块评分（n = 23 个匹配基因，涵盖 \textit{FAP}、\textit{ACTA2}、\textit{COL1A1}、\textit{COL3A1}、\textit{SPP1}、\textit{MMP2}、\textit{LUM}、\textit{DCN} 等 CAF/ECM 重塑相关分子）和两个模块的 Z-score 标准化后加和的组合 niche 评分在 CHC20 空间坐标上的投影结果。三张图呈现出的空间格局高度相似，高评分区域（黄色/橙色）均集中于切片右上方沿组织边缘的同一条带状区域，且低评分区域（深色）均集中于切片下方大片区域，深色区域与 CHC20 的 H\&E 染色切片图肿瘤实质密集区相对应，说明通过 HCC4R 切片筛选出的签名基因能够有效地作为标准识别其他切片不同的生态位。

图 \ref{fig:score_comparison} 小提琴图进一步以组合评分为指标，比较了 HCC4R 与 CHC20 两切片间的整体分布。两组组合 niche 评分分布的 Mann-Whitney U 检验 p 值为 0.094，未达到常规 0.05 显著性水平，两小提琴的中位数（HCC4R 约 −0.5，CHC20 约 −0.2）、四分位距及整体形态均高度重叠，两者的上下边界（HCC4R 约 −2.9 至 5.1，CHC20 约 −4.5 至 7.4）虽然 CHC20 尾部略宽，但主体分布区间基本一致。

综上，CHC20 的高分区域与其组织学切片的肿瘤区域边缘免疫抑制区域在空间上高度重叠，支持所识别的 niche 是 HCC 的普遍性组织学特征而非单一切片特异现象~\citep{llovet2021hepatocellular}。值得关注的是，CHC20 为混合型肝细胞-胆管癌，其免疫微环境组成可能与经典 HCC（HCC4R）存在生物学差异，这一背景下两切片间的评分一致性尤为意义重大。

\subsection{TCGA-LIHC 队列的独立预后验证}

\begin{figure}[h] % [htbp] 允许 LaTeX 自动选择最佳位置（当前位置、顶部、底部或单独一页）
    \centering
    % 第一行左图 (a)
    \begin{subfigure}[b]{0.48\textwidth}
        \centering
        \includegraphics[width=\textwidth]{plots/km_plot.png}
        \caption{ssGSEA 高分组与低分组 KM 曲线}
        \label{fig:km_plot}
    \end{subfigure}
    \hfill % 水平填充，自动推开左右两张图
    % 第一行右图 (b)
    \begin{subfigure}[b]{0.48\textwidth}
        \centering
        \includegraphics[width=\textwidth]{plots/cox_forest_plot.png}
        \caption{ssGSEA 高分组与低分组Cox 森林图}
        \label{fig:cox_forest}
    \end{subfigure}
    \caption{TCGA-LIHC 队列生存分析结果}
    \label{fig:survival_result}
\end{figure}

将 HCC4R 空间签名基因集通过 ssGSEA 投影至 TCGA-LIHC 队列患者，
并依据 ssGSEA 评分中位数将患者分为 High 与 Low 两组。
Kaplan-Meier 生存分析（图 \ref{fig:km_plot}）显示，两组总体生存曲线整体接近，
log-rank 检验未达到统计学显著性（$p = 0.79$），提示该空间签名评分在当前 TCGA-LIHC 队列中未能有效区分患者总体生存差异。

进一步进行多变量 Cox 比例风险回归分析（图 \ref{fig:cox_forest}），
在纳入年龄和肿瘤分期等临床因素后，ssGSEA 评分对应的风险比接近 1
（HR = 0.92，95\% CI: 0.36--2.36，$p = 0.862$），
其置信区间跨越 1，说明该评分在当前模型中未表现出独立预后价值。
相比之下，部分肿瘤分期变量在 Cox 模型中显示出较明显的风险升高趋势，
提示临床分期仍是 TCGA-LIHC 队列中影响总体生存的重要因素。

这个结果说明空间生态位 signature 虽然能反映组织成分，但在 TCGA-LIHC bulk OS 中不是独立预后因子，与生存关系不大。

\section{讨论}

本研究建立了一套贯穿"单细胞参考构建--空间反卷积--生态位识别--机制解析--跨样本与临床验证"五个阶段的完整分析框架，实现了从单细胞转录组参考图谱提取空间免疫调控规律、并在独立组织与患者队列中进行系统性验证的研究目标。与传统以先验标记基因集定义肿瘤免疫抑制状态的分析范式相比，本方案具有三个主要创新点：第一，采用无监督的 Cellular Neighborhoods 聚类策略直接从数据本身发现免疫抑制空间结构，而非依赖于预先设定的细胞类型权重或阈值；第二，针对 Visium 平台固有的细胞稀释效应，引入检出率差值 $\Delta$fraction 作为传统 $\log_2$FC 筛选的补偿维度，有效降低了因稀疏表达导致的假阴性风险，使得 Treg 相关标志基因的特征提取更为敏感；第三，以规则化空间域标签体系作为独立的语义参照系，与数据驱动的 Leiden 聚类结果形成方法学层面的交叉印证，显著增强了免疫抑制生态位判定的稳健性与可解释性。该整合策略为单细胞分辨率参考图谱与空间转录组数据的有机结合提供了可重复的方法学范式，其"发现--机制--验证"闭环设计对其他实体瘤的微环境空间分型研究亦有借鉴意义。

尽管在组织学层面与跨空间切片验证中均获得了高度一致的结果，本研究的 TCGA-LIHC 独立预后验证环节与预期不符：Kaplan-Meier 生存分析（$p = 0.79$）与多变量 Cox 回归（$\text{HR} = 0.92$，$p = 0.862$）均未达到统计学显著性，森林图与 KM 曲线均未能从本空间签名中提取出具有独立预后意义的风险分层信号。这一矛盾可能归因于多个方面的方法学局限。其一，ssGSEA 投影至 bulk 转录组的过程可能会严重稀释空间特异性的信号，bulk RNA-seq 将整张组织匀浆后测量平均表达水平，丢失了空间定位信息，使得真正起驱动作用的局部高信号区域（如 Treg 富集的侵袭前沿或纤维化间质）可能被大面积的正常肝实质或肿瘤实质信号所掩盖~\citep{freije2012gene,gao2020loss}。其二，TCGA-LIHC 队列中的肿瘤组织往往伴有不同程度的坏死、炎症与基质成分比例患者间差异极大，而 bulk 水平的免疫微环境组成受取样位置影响显著~\citep{zhang2020prognostic}；即使同一例患者在影像引导下取两个不同位置的标本，其肿瘤内免疫细胞比例也可能有很大差异，这一取样偏差可能掩盖 bulk ssGSEA 评分与真实总体生存之间的弱效应信号。其三，从空间转录组特征到总体生存的因果链路上存在多个中间环节，免疫抑制生态位可能是肿瘤切除后短期内复发或治疗抵抗的局部驱动因子，但未必能对总体生存造成显著影响~\citep{zhang2020prognostic}。另外，TCGA-LIHC 队列的临床资料多源于不同纳入中心，手术质量、辅助治疗方案及随访完整度存在系统异质性，进一步引入了非生物学层面的随机噪声~\citep{liu2018tcga}。综上，空间 niche signature 虽能精准刻画且高度可复现免疫空间结构，但从其到 bulk OS 预后因子的临床转化可能需要更精细的终点指标（如无复发生存期 RFS）或更均衡的前瞻性样本采集策略~\citep{liu2018tcga}。

除生存分析不显著之外，本章研究在方法学层面还存在以下若干有待改进之处：第一，Immunosuppressive niche score 中度量依赖邻域参数 $k$ 与半径倍数的选择，虽已通过敏感性分析证明结果的总体稳健性（$\rho > 0.90$, $\text{WJ} > 0.80$），但其对极端参数变化的局部空间边界影响仍不可完全忽视，未来研究可考虑利用贝叶斯层次模型同时对后验空间模式与邻域概率分布进行概率推断，降低对单一固定邻域定义的依赖~\citep{kleshchevnikov2022cell2location,rubin2017bayesian}；第二，配体-受体通讯分析虽然引入了秩归一化方法规避了原始乘积法的零值膨胀缺陷，但未能完全区分真正双向通讯与配体--受体基因在同一 shout--spot 内因不同细胞类型各自表达而产生的高共表达假象~\citep{cang2023commot}，可能需要结合配体--受体基因共定位概率模型加以解决；第三，免疫抑制基因模块仅依据侧上调进行筛选，未能顾及抑制性分子下调（如 HLA 基因丢失）等免疫逃逸机制~\citep{montes2021spatial}；第四，现有分析仅纳入了一个 HCC 切片 HCC4R 和一个混合型肝细胞--胆管癌 CHC20 进行跨样本验证，样本量仍显不足，未来工作应扩展至多中心的多癌种空间转录组队列。

展望未来，本研究所建立的分析框架可在以下三个方向加以拓展。第一，通过与多重免疫荧光（mIF）或 MIBI/CODEX 等空间蛋白质组学数据的配准融合，可在蛋白验证层面进一步确认免疫抑制生态位中关键分子（CXCL12--CXCR4 轴、PD-1/PD-L1 免疫检查点）的空间共定位关系，为免疫检查点阻断（ICB）治疗的个体化疗效预测提供分子标志物~\citep{keren2018structured,zhu2021spark}。第二，随着 10x Visium HD 以及 Slide-seq、Stereo-seq 等更高分辨率空间测序技术的日趋成熟，未来可在接近单细胞甚至亚细胞精度上重新评估本研究的结构域定义方案及特征基因签名是否适用于亚 spot 的精细微环境区隔，进一步探索肿瘤--免疫交界面的精细拓扑结构~\citep{moses2022museum,chen2022spatia}。第三，可将"空间生态位 signature"作为输入特征训练 Transformer 等深度学习模型用于手术切缘患者复发风险的无标记预测，或结合数字病理学将宏观 H\&E 图像特征与微观空间转录组 signal 进行多模态整合，构建"宏观--微观"跨尺度的肿瘤免疫状态评估体系~\citep{chen2021liver,shmatko2022artific}。本研究积累的数据处理流程、特征筛选准则以及敏感分析框架可为上述拓展方向提供基础性的方法学参照。

% ============================================================
%  第五部分（附录）：代码实现细节
% ============================================================

\newpage
\appendix
\renewcommand{\thesection}{\Alph{section}}

\section{配体-受体通讯对列表}
\label{app:lr_table}

\begin{table}[H]
\centering
\caption{纳入分析的配体-受体通讯对}
\label{tab:lr_pairs}
\begin{tabularx}{\textwidth}{lllX}
\toprule
\textbf{配体} & \textbf{受体} & \textbf{通讯对} & \textbf{生物学意义} \\
\midrule
CXCL12 & CXCR4   & CXCL12--CXCR4  & CAF 分泌，招募 Treg/MDSC~\citep{domanska2013targeting} \\
TGFB1  & TGFBR1  & TGFB1--TGFBR1  & 免疫抑制细胞因子，CAF+Treg 共分泌~\citep{mariathasan2018tgfb} \\
PDCD1  & CD274   & PD-1--PD-L1    & 经典免疫检查点轴~\citep{sharma2023immunosuppressive} \\
TIGIT  & NECTIN2 & TIGIT--NECTIN2 & Treg 抑制性检查点~\citep{sakaguchi2020regulatory} \\
LAG3   & HLA-DRA & LAG3--MHC-II   & T 细胞耗竭相关~\citep{sharma2023immunosuppressive} \\
CCL22  & CCR4    & CCL22--CCR4    & TAM 招募 Treg~\citep{curiel2004specific} \\
IL10   & IL10RA  & IL10--IL10RA   & Treg 免疫抑制细胞因子轴~\citep{noy2014tumor} \\
SPP1   & CD44    & SPP1--CD44     & TAM 分泌骨桥蛋白~\citep{biswas2012macrophage} \\
MIF    & CD74    & MIF--CD74      & 巨噬细胞迁移抑制因子~\citep{noy2014tumor} \\
VEGFA  & KDR     & VEGFA--KDR     & 血管生成，CAF 相关~\citep{kalluri2016biology} \\
LGALS9 & HAVCR2  & Galectin9--TIM-3 & 抑制性检查点~\citep{sharma2023immunosuppressive} \\
CCL2   & CCR2    & CCL2--CCR2     & 巨噬细胞招募通路~\citep{biswas2012macrophage} \\
\bottomrule
\end{tabularx}
\end{table}

\section{邻域参数敏感性分析指标定义}
\label{app:sensitivity}

正文 2.4 节所述邻域参数敏感性分析采用两个互补的连续型指标，联合评估 $k$-NN 邻居数与半径倍增系数扰动前后综合 niche 评分的一致性。

指标一为 Spearman 秩相关系数 $\rho$，定义为
\begin{equation}
\rho = 1 - \frac{6 \sum_i d_i^2}{n(n^2-1)}
\label{eq:spearman}
\end{equation}
其中 $d_i$ 为参考评分与扰动后评分在第 $i$ 个 spot 上的秩次之差，$n$ 为 spot 总数。$\rho$ 直接度量全部 spot 相对排序的一致性，不受任何截断阈值影响，取值范围为 $[-1, 1]$，越接近 1 表示参数扰动对 spot 排序几乎没有影响。

指标二为连续型加权 Jaccard 相似度（Weighted Jaccard），定义为
\begin{equation}
\text{WJ}(\mathbf{a}, \mathbf{b}) = \frac{\sum_i \min(a_i, b_i)}{\sum_i \max(a_i, b_i)}
\label{eq:weighted_jaccard}
\end{equation}
其中 $\mathbf{a}$、$\mathbf{b}$ 分别为参考评分向量与扰动后评分向量经 Min-Max 归一化至 $[0,1]$ 后的结果。相较于传统基于二值集合的 Jaccard 指数，该连续形式对高分区域（即候选 niche 区域）天然赋予更高权重，弥补了 Spearman 秩相关对顶部与底部区域同等看待的不足，取值范围为 $[0, 1]$，越接近 1 表示高分区域的空间分布高度重合。

\section{辅助评分公式}
\label{app:scores}

辅助评分体系定义如下，其中 $z(\cdot)$ 表示跨 spot Z-score 标准化：

\begin{align}
S_{\text{Treg-like},i} &= z(\hat{p}_{\text{Treg},i}) + z(\hat{s}_{\text{gene},i}) \\
S_{\text{immune-stroma},i} &= z(\hat{p}_{\text{Treg},i}) + z(\hat{p}_{\text{T/NK},i}) + z(\hat{p}_{\text{Myeloid},i}) + z(\hat{p}_{\text{Fibroblast},i}) \\
S_{\text{niche},i} &= z(\hat{p}_{\text{Treg},i}) + z(\hat{p}_{\text{Myeloid},i}) + z(\hat{p}_{\text{Fibroblast},i}) + z(\hat{s}_{\text{gene},i}) + z(\bar{n}_{\text{Hep},i})
\end{align}

其中 $\hat{p}$ 为 Cell2location 后验丰度归一化比例，$\hat{s}_{\text{gene}}$ 为免疫抑制基因模块评分，
$\bar{n}_{\text{Hep}}$ 为 spot 邻域内肝细胞比例均值（用于捕捉肿瘤-基质界面效应）。
双模块评分（用于 CHC20 验证）：

\begin{align}
\text{immune\_score}_i &= \frac{1}{|G_{\text{immune}}|}\sum_{g \in G_{\text{immune}}} E_{ig} \\
\text{stromal\_score}_i &= \frac{1}{|G_{\text{stromal}}|}\sum_{g \in G_{\text{stromal}}} E_{ig} \\
\text{combined\_niche\_score}_i &= z(\text{immune\_score}_i) + z(\text{stromal\_score}_i)
\end{align}

$E_{ig}$ 为 log1p 标准化（CP10K）表达量；免疫模块21个基因涵盖 \textit{FOXP3}、\textit{IL2RA}、\textit{CTLA4}、\textit{CD163} 等，
基质/ECM 模块21个基因涵盖 \textit{FAP}、\textit{ACTA2}、\textit{COL1A1}、\textit{SPP1} 等。

\section{核心代码实现}
\label{app:code}

\subsection{Cell2location 反卷积}

\begin{lstlisting}[caption={RegressionModel 训练与空间建模核心代码}]
# Stage 1: 训练参考签名矩阵
from cell2location.models import RegressionModel
RegressionModel.setup_anndata(adata_sc, labels_key='final_celltype')
mod = RegressionModel(adata_sc)
mod.train(max_epochs=250, batch_size=1024, use_gpu=True,
          early_stopping=True, early_stopping_patience=30)
cell_state_df = mod.export_posterior(adata_sc,
    use_cols=mod.feature_in_columns)['means_per_cluster_mu_fg']

# Stage 2: 空间建模
import cell2location
cell2location.models.Cell2location.setup_anndata(adata_vis)
mod = cell2location.models.Cell2location(adata_vis,
    cell_state_df=cell_state_df,
    N_cells_per_location=30, detection_alpha=20)
mod.train(max_epochs=1000, batch_size=1000, use_gpu=True)
adata_vis = mod.export_posterior(adata_vis,
    sample_kwargs={'num_samples': 1000, 'batch_size': 1000})
# 结果存储在 adata_vis.obsm['means_cell_abundance_w_sf']
\end{lstlisting}

\subsection{邻域组成聚类与生态位注释}

\begin{lstlisting}[caption={邻域组成向量计算、Leiden 聚类与综合排名注释}]
from scipy.spatial import KDTree
import scanpy as sc

# 构建 kNN 邻接图并计算邻域组成向量
tree = KDTree(coords)
_, indices = tree.query(coords, k=k + 1)  # 排除自身
neighbors = [indices[i, 1:] for i in range(len(coords))]
arr = proportions.to_numpy()
neighborhood_comp = np.array(
    [arr[idx].mean(axis=0) if len(idx) else arr[i]
     for i, idx in enumerate(neighbors)])

# Leiden 聚类（直接在组成向量特征空间，无需 PCA）
adata_nc = ad.AnnData(neighborhood_comp.astype(float))
sc.pp.neighbors(adata_nc, n_neighbors=n_neighbors, use_rep='X')
sc.tl.leiden(adata_nc, resolution=0.5, key_added='nc_cluster')

# 综合排名法注释免疫抑制 niche
stats = df.groupby('nc_cluster')[
    [treg_col, myeloid_col, fibroblast_col, gene_score_col]].mean()
for col in stats.columns:
    stats[f'rank_{col}'] = stats[col].rank(ascending=False)
stats['combined_rank'] = stats.filter(like='rank_').sum(axis=1)
niche_id = stats['combined_rank'].idxmin()
\end{lstlisting}

\subsection{三层特征基因筛选}

\begin{lstlisting}[caption={Layer 1 差异检验、Layer 3 Gini Index 计算}]
# Layer 1: block-level Mann-Whitney U + AUC
from scipy.stats import mannwhitneyu
# 空间网格聚合（降低空间自相关）
gx = np.floor(df['spatial_x'].values / block_size).astype(int)
gy = np.floor(df['spatial_y'].values / block_size).astype(int)
block_ids = [f'{x}_{y}' for x, y in zip(gx, gy)]
# 计算 AUC = U / (n_high * n_low)
u_stat, pval = mannwhitneyu(high_vals, low_vals, alternative='greater')
auc = u_stat / (len(high_vals) * len(low_vals))

# Layer 3: Gini Index 筛选局灶性稀疏基因
def gini_index(x):
    x = np.sort(x.ravel())
    n = len(x)
    cumx = np.cumsum(x)
    return (n + 1 - 2 * cumx.sum() / cumx[-1]) / n if cumx[-1] > 0 else 0
\end{lstlisting}

\subsection{双模块评分与 ssGSEA}

\begin{lstlisting}[caption={双模块评分计算（Python）与 ssGSEA 评分（R）}]
# Python: 双模块评分
def compute_dual_module_scores(adata, immune_genes, stromal_genes):
    expr = _expression_frame(adata)  # CP10K + log1p
    immune_score = expr[[g for g in immune_genes
                          if g in expr.columns]].mean(axis=1)
    stromal_score = expr[[g for g in stromal_genes
                           if g in expr.columns]].mean(axis=1)
    def _z(s): return (s - s.mean()) / s.std() if s.std() > 0 else s * 0
    return immune_score, stromal_score, _z(immune_score) + _z(stromal_score)
\end{lstlisting}

\begin{lstlisting}[language=R, caption={ssGSEA + 多变量 Cox 回归（R）}]
library(GSVA); library(survival); library(survminer)
gsva_scores <- gsva(expr_matrix, list(sig = signature_genes),
    method='ssgsea', kcdf='Gaussian', abs.ranking=TRUE, min.sz=5)
cox_model <- coxph(
    Surv(OS_time, OS_event) ~ ssgsea_score + age + stage,
    data = clinical_data)
summary(cox_model)
\end{lstlisting}

% ============================================================
% 参考文献
% ============================================================
\bibliographystyle{unsrt}
\bibliography{report_0701}

\end{document}
